\section{The 5-State Epistemic Assessment Model}
\label{sec:assessment}

Rather than flattening heterogeneous findings into an arbitrary compliance percentage, Reg2App evaluates every statutory requirement against a discrete, 5-state epistemic assessment space $\mathcal{S} = \{\mathbf{S}, \mathbf{P}, \mathbf{I}, \mathbf{O}, \mathbf{A}\}$~\cite{fagin_epistemic_logic, belnap_four_valued}, grounded in the boundaries of client-side program observability.

\subsection{Epistemic State Space and Formal Semantics}
\label{sec:assessment:states}

For each requirement $R_i$ evaluated against application $A$, the assessment function emits an epistemic state $s(R_i, A) \in \mathcal{S}$:
\begin{equation}
\label{eq:state:semantics}
\small
s(R_i, A) = \begin{cases}
\mathbf{A}, & \text{if } \text{Sector}(A) \notin \text{Scope}(R_i) \\
\mathbf{O}, & \text{if } \Omega(R_i).L_i = \mathbf{NONE} \\
\mathbf{P}, & \text{if } \exists w \in W(A, R_i) \text{ s.t. } \neg \phi_{P_i}(w) \\
\mathbf{I}, & \text{if } (\exists w, \mathcal{D}(\phi_{P_i}, w) = \bot) \vee (W = \emptyset \wedge \text{Aff}(R_i)) \\
\mathbf{S}, & \text{if } \forall w, \phi_{P_i}(w) \wedge (W \neq \emptyset \vee \text{Neg}(R_i))
\end{cases}
\end{equation}
Equation~(\ref{eq:state:semantics}) is evaluated in listed priority order (the first matching condition determines the emitted state). Here, $W(A, R_i) \subseteq V_{\text{Witness}}$ is the set of extracted client witnesses, and $\mathcal{D}: \mathcal{P} \times V_{\text{Witness}} \rightarrow \{\text{True}, \text{False}, \bot\}$ evaluates predicate $\phi_{P_i}$ over witness $w$. Under these operational semantics:
(1)~\textbf{Potential Non-Conformance} ($\mathbf{P}$) takes precedence: if at least one observable witness violates the technical safety property ($\exists w, \neg \phi(w)$), the clause is flagged as non-conforming. Because safety properties demand universal compliance across all call sites (e.g., all network endpoints must enforce encryption), a single non-conforming witness immediately refutes compliance. Evaluating $\mathbf{P}$ prior to $\mathbf{I}$ ensures that observable defects are never masked by undecidable reflection stubs elsewhere in the application.
(2)~\textbf{Insufficient Evidence} ($\mathbf{I}$) subsumes two operational causes: (i)~\textit{analysis undecidability}, where code is present but dynamic reflection, class loading, or native C/C++ stubs~\cite{octeau_reflection, ernst_static_dynamic} force $\mathcal{D} = \bot$; and (ii)~\textit{feature absence}, where affirmative obligations ($W = \emptyset \wedge \text{Aff}(R_i)$, e.g., biometric prompt or secure database) are not instantiated client-side. Because static analysis cannot distinguish uninstantiated features from obfuscated ones without speculation, both yield $\mathbf{I}$, signaling unverified paths. Consequently, Evidence Coverage ($EC$) quantifies the proportion of rules yielding a definitive verdict ($\mathbf{S}$ or $\mathbf{P}$), partially capturing client feature presence.
(3)~\textbf{Supported} ($\mathbf{S}$) requires affirmative client evidence and universal compliance across all evaluated call sites ($\forall w, \phi(w) = \text{True}$). For negative prohibitions (e.g., banning backup extraction, cleartext HTTP, or deprecated ciphers), the absence of insecure invocations under secure configuration validates compliance ($W \neq \emptyset \vee \text{IsNegativeProhibition}(P_i)$). For requirements tagged as $\mathbf{PARTIAL}$, $\mathbf{S}$ denotes \textit{bounded client-side conformance} ($\mathbf{S}_{\text{client}}$): it certifies that all observable client preconditions are fulfilled, while preserving the epistemic boundary that backend server or organizational obligations remain unverified off-device. Thus, $\mathbf{S}$ and $\mathbf{P}$ are strictly mutually exclusive.
(4)~\textbf{Not Observable} ($\mathbf{O}$) signifies that the clause requires backend, organizational, or physical evidence ($L_i = \mathbf{NONE}$), rendering client-side verification impossible.
(5)~\textbf{Not Applicable} ($\mathbf{A}$) indicates that the requirement is sector-scoped (e.g., central bank mandates) and the target application belongs to an exempt category.

Formally, the states define an information preorder $(\mathcal{S}, \sqsubseteq)$ in domain theory~\cite{davey_priestley_lattices, belnap_four_valued}: $\mathbf{InsufficientEvidence}$ ($\mathbf{I}$) represents the epistemic bottom $\bot$ of incomplete client knowledge ($\mathbf{I} \sqsubseteq \mathbf{S}$ and $\mathbf{I} \sqsubseteq \mathbf{P}$), while $\mathbf{Supported}$ ($\mathbf{S}$) and $\mathbf{PotentialNonConformance}$ ($\mathbf{P}$) are maximal, mutually exclusive informative states. States $\mathbf{O}$ and $\mathbf{A}$ represent scope boundaries, excluded from client verification denominators.

\subsection{Formal Assessment Metrics}
\label{sec:assessment:metrics}

Let $N_{\text{total}} = N_S + N_P + N_I + N_O + N_A$ denote total clauses in the active knowledge base. The set of relevant, client-observable requirements is:
\begin{equation}
\label{eq:assessment:obs}
N_{\text{Observable}} = N_S + N_P + N_I
\end{equation}
where $N_O$ (backend/legal clauses) and $N_A$ (out of sector scope) are strictly excluded from client verification denominators.

We define three objective, provenance-grounded metrics:

\paragraph{1) Evidence Coverage ($EC$)} Quantifies verification completeness over observable requirements:
\begin{equation}
\label{eq:assessment:ec}
EC = \frac{N_S + N_P}{N_S + N_P + N_I} = \frac{N_S + N_P}{N_{\text{Observable}}} \in [0, 1]
\end{equation}
When $EC = 1.00$, the engine encountered zero unresolvable barriers across observable requirements.

\paragraph{2) Evidence-Supported Compliance ($ESC$)} Measures the proportion of verified requirements demonstrating affirmative compliance:
\begin{equation}
\label{eq:assessment:esc}
ESC = \frac{N_S}{N_S + N_P} \in [0, 1] \quad (\text{if } N_S + N_P = 0, ESC = 0.0)
\end{equation}
Crucially, $ESC$ depends exclusively on verifiable client-side evidence, strictly excluding unobservable backend clauses ($N_O$) and unresolvable code paths ($N_I$). While $ESC$ can be theoretically stratified into $ESC_{\text{FULL}} = \frac{N_{S|\text{FULL}}}{N_{S|\text{FULL}} + N_{P|\text{FULL}}}$ (7 rules) and $ESC_{\text{PARTIAL}} = \frac{N_{S|\text{PARTIAL}}}{N_{S|\text{PARTIAL}} + N_{P|\text{PARTIAL}}}$ (5 rules), our empirical evaluations report pooled $ESC$ across all active rules for cross-sectoral comparability.

\paragraph{3) Strict Conformance Index ($SCI$)} Provides a conservative lower-bound metric, penalizing incomplete evidence ($N_I$) as non-compliant:
\begin{equation}
\label{eq:assessment:sci}
SCI = \frac{N_S}{N_S + N_P + N_I} \in [0, 1]
\end{equation}
Always $SCI \le ESC$, with equality if and only if $N_I = 0$ ($EC = 1.0$). The divergence $\Delta_{\text{epistemic}} = ESC - SCI = \frac{N_S \cdot N_I}{(N_S + N_P)(N_S + N_P + N_I)}$ quantifies the epistemic deficit from unresolvable constructs.

\paragraph{Confidence Intervals and Interpretation} We compute two-sided $95\%$ confidence intervals via the Wilson score interval~\cite{wilson_score_interval}:
\begin{equation}
\text{CI}_{95\%}(\hat{p}) = \frac{\hat{p} + \frac{z^2}{2n} \pm z \sqrt{\frac{\hat{p}(1-\hat{p})}{n} + \frac{z^2}{4n^2}}}{1 + \frac{z^2}{n}} \quad (z = 1.960)
\end{equation}
We define three operational thresholds: \textit{Adequate Observability} ($EC \ge 0.80$), confirming static analysis resolved the vast majority of client logic; \textit{High Epistemic Uncertainty} ($EC < 0.70$), signaling that supplemental dynamic analysis is required; and \textit{High Technical Conformance} ($ESC \ge 0.85$ with $EC \ge 0.80$).

\subsection{Architectural Invariants and Properties}
\label{sec:assessment:invariants}

\noindent\textbf{System Invariant 1 (Epistemic Anti-Hallucination Guard).}
\textit{For any application $A$ and statutory corpus $\mathcal{K}$, the assessment engine is guaranteed never to attribute affirmative compliance ($\mathbf{S}$) or non-conformance ($\mathbf{P}$) to a requirement residing beyond client-side observational boundaries ($L_i = \mathbf{NONE}$):
\begin{equation}
\forall R_i \text{ s.t. } \Omega(R_i).L_i = \mathbf{NONE} \implies s(R_i, A) \equiv \mathbf{O}
\end{equation}}
Enforced by an architectural guard check in Algorithm~\ref{alg:assessment} prior to analyzer invocation, this invariant eliminates compliance hallucinations by construction.

\noindent\textbf{Metric Property 2 (Monotonicity of Evidence Coverage).}
\textit{Let $\mathcal{A}_1 \subseteq \mathcal{A}_2$ be two static engines where $\mathcal{A}_2$ resolves a strict superset of previously undecidable constructs. Then $EC(\mathcal{A}_1) \le EC(\mathcal{A}_2)$, confirming that $EC$ monotonically reflects analyzer resolving power.}

\subsection{Assessment Algorithm}
\label{sec:assessment:algorithm}

Algorithm~\ref{alg:assessment} outlines the evaluation procedure. The engine consults observability metadata first, classifying $L_i = \mathbf{NONE}$ as $\mathbf{NotObservable}$ without analyzer invocation (Invariant~1). For observable clauses, analyzers collect evidence witnesses and prioritize defect detection to enforce mutual exclusivity.

\begin{figure}[t]
\begin{tcolorbox}[colback=gray!5!white,colframe=black!75,title={\textbf{Algorithm 1}: Epistemic Assessment Procedure},fontupper=\scriptsize]
\textbf{Input:} Application $A$, Knowledge Base $\mathcal{K}$, Analyzers $\mathcal{A}$\\
\textbf{Output:} Assessment vector $\vec{s}(A)$, Graph $G$
\begin{algorithmic}
\FOR{each clause $R_i \in \mathcal{K}$}
  \IF{$\text{AppSector}(A) \notin \text{Scope}(R_i)$}
    \STATE $s_i \leftarrow \mathbf{NotApplicable}$
  \ELSIF{$\Omega(R_i).L_i = \mathbf{NONE}$}
    \STATE $s_i \leftarrow \mathbf{NotObservable}$
  \ELSE
    \STATE $W \leftarrow \text{Analyze}(\mathcal{A}, A, R_i)$ \COMMENT{collect witnesses}
    \IF{$\exists w \in W \text{ s.t. } \phi_{P_i}(w) = \text{False}$}
      \STATE $s_i \leftarrow \mathbf{PotentialNonConformance}$
    \ELSIF{$\exists w \in W, \mathcal{D}(\phi_{P_i}, w) = \bot \text{ \OR } (W = \emptyset \wedge \text{IsAffirmative}(P_i))$}
      \STATE $s_i \leftarrow \mathbf{InsufficientEvidence}$
    \ELSE
      \STATE $s_i \leftarrow \mathbf{Supported}$
    \ENDIF
    \STATE $G \leftarrow G \cup \{(v_s^i, v_c^i, v_p^i, W)\}$
  \ENDIF
\ENDFOR
\STATE \textbf{return} $\vec{s}(A) = (s_1, \ldots, s_n)$, $G$
\end{algorithmic}
\end{tcolorbox}
\refstepcounter{algorithm}\label{alg:assessment}
\captionof{figure}{\textbf{Algorithm \thealgorithm}: Epistemic assessment procedure, showing the observability guard and evidence provenance graph construction.}
\end{figure}
